% `tex.web` §962 folds every pattern character through `\lccode`:
%
%     if cur_chr="." then cur_chr:=0 {edge-of-word delimiter}
%     else begin cur_chr:=lc_code(cur_chr);
%       if cur_chr=0 then
%         begin print_err("Nonletter");
%         help1("(See Appendix H.)"); error;
%         end;
%       end;
%
% and §963 keeps the *last* of two patterns for the same letters, saying so:
%
%     if trie_o[q]<>min_quarterword then
%       begin print_err("Duplicate pattern");
%       help1("(See Appendix H.)"); error;
%       end;
%     trie_o[q]:=v;
%
% The two rules are one bug together.  Without the folding, `bb1` and `0B2B0`
% are two patterns where TeX has one — so TeX keeps the second, whose values
% are all even, and offers no break at all where we offered four.  The
% overwrite itself we already had, being what a map insert does; what was
% missing was the key it happens on.
%
% A pattern character with no `\lccode` is diagnosed and kept as character
% zero, with `Nonletter` said once for each. §960 also rejects commands other than
% letter/other-character tokens, and §962 treats a second consecutive digit as
% a letter; its zero `\lccode` therefore raises `Nonletter`.
\catcode`\{=1 \catcode`\}=2
\font\r=cmr10 \r \hsize20pt \parindent0pt \baselineskip12pt
\tolerance10000 \pretolerance-1 \hbadness10000 \uchyph=1
\showboxbreadth99 \showboxdepth99 \tracingonline0
\patterns{bb bb1 0B2B0 b1c}
\setbox0\vbox{\noindent x BBBBBB\par}\showbox0
\lccode`\!=`\!
\patterns{a1!b !1a}
\patterns{x1y x2y}
\chardef\badchar=`b
\patterns{\badchar 50b}
\end
